\documentclass{exam}
\usepackage{mystyle}
\usepackage{placeins}

\title{MATH 3191 Written Homework 5\\\rule{30pt}{1pt}/20 points}
\author{Name: \rule{150pt}{1pt}}
\date{Due: October 11th}

\begin{document}
\maketitle
Show all work that leads to your answers. I will be providing feedback on the work as well as the answers!
\begin{questions}
    \question (20 points total):
    Polynomial Interpolation\footnote{\url{https://en.wikipedia.org/wiki/Polynomial_interpolation}}
    is a useful tool, so we will look at how to use different bases before we get to it with
    Python Lab 4!

    The main idea behind polynomial interpolation is that we want a polynomial to go through a given
    set of points when we either don't necessarily know what the underlying function is or don't care
    about the function itself, and instead only need an approximation.

    For both of the following parts, we will be dealing with the following vector space:
    $V = \mathcal{P}_2(\mathbb{R}), F=\mathbb{R}$ with the standard operations, and the standard
    basis vectors are: $\mathcal{S} = \setBasic{1, t, t^2}$. We will use the following points
    that we want to interpolate: (Note you do not need to use these numbers for anything in the problem,
    but I am providing them so you can come back and see how this was made once you finish Python Lab 4)

    \begin{center}
        \begin{tabular}{|c|c|}
            \hline
            $x$  & $y$ \\
            \hline
            \hline
            $-2$ & $5$ \\
            $-1$ & $2$ \\
            $1$  & $7$ \\
            \hline
        \end{tabular}
    \end{center}
    \begin{parts}
        \part[10]
        The Newton basis for the above interpolation points are given by:

        \[
            \mathcal{N} = \setBasic{1, t+2, t^2 +3t + 2}
        \]

        Similarly, the Lagrange basis is given by:

        \[
            \mathcal{L} = \setBasic{\frac{5}{3}\left(t^2-1\right), -t^2-t+2, \frac{7}{6}\left(t^2+3t+2\right)}
        \]

        Demonstrate that both of these sets $\mathcal{N}$ and $\mathcal{L}$ are actually bases of
        $\mathcal{P}_2(\mathbb{R})$
        \vfill
        \newpage
        \part[10]
        Construct $P(\mathcal{N})_\mathcal{L} = \underset{\mathcal{L}\leftarrow\mathcal{N}}{P}$
        \vfill
    \end{parts}
\end{questions}
\end{document}
